\begin{statementbox}
	\textbf{\textcolor{CStatement}{条件}}
	\begin{description}[style=nextline, leftmargin=2.8em, labelsep=0.5em, itemsep=0.35em]
		\item[\textbf{\textcolor{CStatement}{条件 1（正方形45度）：}}]
			正方形 $ABCD$ 中，$E$ 在边 $BC$ 上，$F$ 在边 $CD$ 上，且 $\angle EAF = 45^\circ$。
		\item[\textbf{\textcolor{CStatement}{条件 2（变量设定）：}}]
			设 $AB = a$，$BE = x$，$DF = y$，$EF = z$。
	\end{description}
	\textbf{\textcolor{CStatement}{结论}}
	\begin{description}[style=nextline, leftmargin=2.8em, labelsep=0.5em, itemsep=0.45em]
		\item[\textbf{\textcolor{CStatement}{结论一（线段和关系）：}}]
			\textbf{\textcolor{CStatement}{直观描述：}}
			点 $E$ 到 $F$ 的线段长等于 $BE$ 与 $DF$ 之和。
			\[
				EF = BE + DF
			\]
		\item[\textbf{\textcolor{CStatement}{结论二（乘积关系）：}}]
			\textbf{\textcolor{CStatement}{直观描述：}}
			两线段 $x$ 与 $y$ 满足二次方程。
			\[
				xy = a^2 - a(x+y)
			\]
		\item[\textbf{\textcolor{CStatement}{推广结论（变形乘积关系）：}}]
			\textbf{\textcolor{CStatement}{直观描述：}}
			$x$ 与 $y$ 分别加边长 $a$ 后的乘积为定值 $2a^2$。
			\[
				(x+a)(y+a) = 2a^2
			\]
	\end{description}

\end{statementbox}
